\Needspace{10\baselineskip}
\begin{center}
PALLAS.
\end{center}


Monatliche Correspondenz. Vol. VII, pp.\ 369—375 and 466—469. April and May 1803.


[The April issue gives the substance of the letter to \textsc{Bode} of 3 March. See above, p.\ 233. The May issue also contains the following passage:]

— — — These observations [by Dr.\ \textsc{Olbers}] still agree remarkably well with Dr.\ \textsc{Gauss}’s elements VI for the orbit. Here is the comparison:

\begin{center}
\begingroup\fontsize{10}{12}\selectfont\setlength{\tabcolsep}{3pt}\renewcommand{\arraystretch}{1.12}
\begin{tabular}{lrl}
\multicolumn{3}{l}{1803 March 4. $17^u\,11^m\,41^s$ Bremen mean time} \\
Calculated R.A. $\PallasSymbol$ & $275^\circ\,52'\,35''9$ & Difference. R—B, $-2''1$ \\
Calculated decl. $\PallasSymbol$ & $8\ 58\ 18.1$ & Difference. R—B, $-4.9$ \\
\end{tabular}\endgroup
\end{center}

Since the error in the earlier elements V is so small and is also always known with sufficient accuracy, the ephemeris given in our December issue will be quite sufficient for finding the planet throughout the period it covers. The right ascension in it is now $2\tfrac{1}{2}$ min.\ too large, and on 9 August it will be $4'\,6''$ too large, insofar as one can rely on the continued agreement of the new elements. — — — Dr.\ \textsc{Gauss} has meanwhile extended the ephemeris of \emph{Pallas} [in the form he has used so far], using the new elements VI [from 9 Aug.] to 23 Oct.\ 1803. One will probably be unable to follow it for longer, since as early as 18 Sept.\ it will begin to be fainter than it was when it was rediscovered.

\begin{flushright}
[\textsc{von Zach}.]
\end{flushright}

